% ECONOMICS_PROBLEM: {"id": "C73-51", "title": "Uniform coalition reachability with unknown population size", "jel_code": "C73", "source_id": "OP-0170"}
% !TeX program = xelatex
\documentclass[11pt,a4paper]{article}
\usepackage[margin=23mm]{geometry}
\usepackage{fontspec}
\setmainfont{Latin Modern Roman}
\usepackage{amsmath,amssymb,mathtools}
\usepackage{xcolor,enumitem,booktabs,longtable,array,xurl,hyperref}
\definecolor{muted}{HTML}{53636C}
\hypersetup{colorlinks=true,urlcolor=blue,linkcolor=blue}
\setlength{\parindent}{0pt}
\setlength{\parskip}{5pt}
\newcommand{\R}{\mathbb R}
\newcommand{\E}{\mathbb E}
\newcommand{\N}{\mathbb N}
\newcommand{\ind}{\mathbf 1}
\newcommand{\Var}{\operatorname{Var}}
\newcommand{\Cov}{\operatorname{Cov}}
\newcommand{\supp}{\operatorname{supp}}
\DeclareMathOperator*{\argmin}{arg\,min}
\DeclareMathOperator*{\argmax}{arg\,max}
\newcommand{\entrymeta}[3]{{\sffamily\small\color{muted}\textbf{\texttt{#1}}\quad|\quad #2\par #3\par}}
\newcommand{\Problemref}[1]{\texttt{#1}}
\newcommand{\JELref}[2]{\texttt{#2} (JEL \texttt{#1})}
\begin{document}
\section*{Uniform coalition reachability with unknown population size}
\textbf{Problem ID:} \texttt{C73-51}\quad \textbf{JEL:} \texttt{C73}\par
% BEGIN ECONOMICS STATEMENT
\label{problem:OP-0170}
% GAME_THEORY_PROBLEM: {"local_id": "GT-040", "original_number": 40, "kind": "P", "entry_kind": "statement_only", "published": false, "review_required": true, "category": "game_theory"}
\label{game:GT-040}
{\sffamily\small\color{muted}
{\sffamily\small\color{muted}\textbf{GT-040} | Parameterized concurrent games\par P | Source-reported general finite-arena decision question\par}
\par}
\subsection*{Definitions and mathematical statement}
An arena has finite vertices $V$, finite action alphabet $\Sigma$, initial vertex $v_0$, target set $F\subseteq V$, and a regular language $L(v,w)\subseteq\Sigma^+$ for every ordered vertex pair. Each language is given by a finite automaton. The arena is complete: for every $v$ and nonempty word $u\in\Sigma^+$ there is a $w$ with $u\in L(v,w)$. If several such $w$ exist, the environment selects one adversarially.

There are $k\geq1$ cooperative agents, indexed $1,\ldots,k$; $k$ is fixed during a play but unknown to the strategy. A uniform coalition strategy is a sequence of functions $\sigma_m:V^+\to\Sigma$, one for every possible agent index $m\geq1$. Equivalently, $\sigma(h)\in\Sigma^{\mathbb N}$ is an infinite action word. For population $k$, only its first $k$ symbols are played. Define
\[
 \operatorname{Out}_k(v_0,\sigma)=
 \{v_0v_1\ldots:\forall t\geq0,\ 
 \sigma_1(v_0\ldots v_t)\cdots\sigma_k(v_0\ldots v_t)
 \in L(v_t,v_{t+1})\}.
\]
The decision question, for the explicit finite-automaton input, is whether
\[
 \exists\sigma\quad\forall k\geq1\quad
 \forall\rho\in\operatorname{Out}_k(v_0,\sigma)\quad
 \exists t\geq0:\rho_t\in F.
\]
The same strategy must work at every fixed population size; its actions can use the agent index and public vertex history, but not $k$.
\subsection*{Short English statement}
Can one coalition strategy force eventual reachability for every unknown population size in an arbitrary finite arena?
\subsection*{Variants and scope boundaries}
The reaching time may depend on the population and outcome; a single uniform time bound is not required. Arbitrary history-dependent strategies are allowed, with no finite-memory promise. The general arena problem differs from the specialized repeated-game construction solved by the source and from one distinguished agent playing against the population.
\subsection*{Source relationship and status}
[S1] distinguishes its solved special model from the still-open arbitrary finite-graph coalition reachability problem. [S2] provides the regular-language arena and population-uniform strategy definitions used to make that general question precise.
\subsection*{References}
\begingroup\footnotesize\raggedright
\textbf{[S1]} Nathalie Bertrand, Patricia Bouyer, Luc Lapointe and Corto Mascle, \emph{Reach Together: How Populations Win Repeated Games}, arXiv:2510.02984 (2025). Locators: Section 3.3, pp. 7--8 (general-arena distinction); Section 6, pp. 15--16 (scope and remaining questions). \url{https://arxiv.org/pdf/2510.02984}

\textbf{[S2]} Nathalie Bertrand, Patricia Bouyer and Anirban Majumdar, \emph{Synthesizing Safe Coalition Strategies}, FSTTCS 2020, LIPIcs 182, Article 39, pp. 39:1--39:17, doi:10.4230/LIPIcs.FSTTCS.2020.39. Locators: Section 2, Definition 1 and the strategy/outcome definitions, pp. 39:3--39:4; Section 4, p. 39:15 (reachability and nonuniform reaching times). \url{https://drops.dagstuhl.de/storage/00lipics/lipics-vol182-fsttcs2020/LIPIcs.FSTTCS.2020.39/LIPIcs.FSTTCS.2020.39.pdf}

% Corrections and retrieval limits for assembly:
% GT-021: external-regret/empirical-mixture statement is editorial; source local
% weighted regret and recursive correlated-policy output are not silently equated.
% GT-022: FP-DQN convergence is not a closed proposition without schedules,
% finite-sample asymptotics and function-approximation/optimizer specifications.
% GT-023: use v5 dated 1 Oct 2026, not v1; Jiaru Li is a v5 author;
% Theorem 5.8 bounds squared SPR, and Theorem 4.5 is approximate representation.
% GT-026: finite-horizon uniform payoff, expected pathwise liminf and liminf
% expected average are separate definitions, not a slash-separated criterion.
% GT-028/029: source strategies observe state histories; SPE includes off-path
% feasible histories and retains previously attained target payoffs.
% GT-030: unrestricted threshold-constrained NE decision is non-r.e. by
% Ummels--Wojtczak 2011, Theorem 4.9; source only motivates unconstrained existence.
% GT-031/032: positional means PURE memoryless; stochastic stationary existence
% is a known larger-strategy result, not the open positional target.
% GT-034: Inria paragraph is an agenda, not a complete computational model.
% GT-035: two-player exact stationary NE is PPAD-complete already; retain >=3
% exact real search as the source's remaining complexity question.
% GT-036: discrete bounded representative supplied; unrestricted continuous-time
% version needs path regularity, stopping-time and simultaneous-stop conventions.
% GT-037/038: tail-measurable maxmin equality does not ensure existence of value.
% GT-039: comparator is the same role in a parallel identical game, not opponent.
% GT-040: k is fixed across a play, unknown to the uniform strategy; no uniform
% reaching time is required. All sources retrieved; later literature not exhaustive.


% Entries 041--060: source-checked formulations; local reference labels restart.
\par\endgroup

\subsection*{Common conventions: Source mathematical conventions}

\textbf{Finite games.} For a finite set $A$, $\Delta(A)$ is its probability
simplex. Players choose independent mixed strategies in Nash equilibrium;
correlation is allowed only when explicitly part of the solution concept.
In a game with payoff functions $u_i$ and strategy profile $x$, an additive
$\varepsilon$ Nash equilibrium satisfies
\[
 u_i(x)\geq u_i(x_i',x_{-i})-\varepsilon
 \quad\text{for every player $i$ and every admissible deviation $x_i'$}.
\]
For numerical approximation, payoffs are normalized as locally specified.
An $\varepsilon$ payoff residual is distinct from distance at most
$\varepsilon$ to the set of exact equilibria. Point convergence, convergence
of empirical play, and convergence in distance to an equilibrium set are
also distinct.

\textbf{Computational representations.} Unless stated otherwise, a complexity
question takes rational data with binary encoding; polynomial time is measured
in total bit length. A fixed-accuracy polynomial algorithm is not necessarily
polynomial in $1/\varepsilon$, and neither is necessarily polynomial in
$\log(1/\varepsilon)$. Succinct games and value, demand, or sampling oracles
must declare their interfaces. Exact real solutions require an explicit
representation; an unspecified list of arbitrary real numbers is not a
finite computational output. A claimed algorithmic barrier is meaningful
only for its specified model and reductions.

\textbf{Dynamic games.} Histories, information, observation, and allowable
randomization are part of the model. A stationary strategy depends only on
the current state; a positional strategy in a deterministic graph game
chooses one action at each controlled vertex. Nash equilibrium at an initial
state and equilibrium after every history are different requirements.
An expectation of a pathwise $\liminf$ payoff, a $\liminf$ of expected averages,
a limiting expected average, and a uniform finite-horizon equilibrium are
different criteria. None is silently substituted for another.

\textbf{Allocation and mechanisms.} A complete allocation assigns every item
exactly once unless otherwise stated. Zero-valued goods, negative values,
capacity constraints, feasibility, lotteries, and copies of goods affect
fairness definitions and are specified locally. Dominant-strategy incentive
compatibility, Bayesian incentive compatibility, universal truthfulness,
and truthfulness in expectation impose different inequalities. Whenever
an objective ratio has zero denominator, its convention is stated or those
instances are excluded.

\textbf{Infinite-dimensional models.} Expectations, integrals, stochastic
equations, and optimization objectives require their local regularity,
integrability, filtrations, and admissible domains. A backward--forward
mean-field system is a boundary-value problem; it is not an initial-value
semiflow without an additional construction. Long-time, large-population,
and vanishing-noise limits require an order of limits and a convergence
topology. If a minimum need not be attained, the objective is its infimum
or supremum and the approximation requirement is stated separately.
% END ECONOMICS STATEMENT
\end{document}
